% ED26-BM-OHTR2012-5: Doumas--Hummel 2012, complete printed 52--66.
\chapter{Computational Models of Higher Cognition}
\label{sec:p19-computational-cognitive-models}

Doumas and Hummel compare traditional symbolic, traditional connectionist,
and symbolic-connectionist process models. Their explanatory target is
the representation and processing of relational information: how a system
represents not just the entities present but which entity occupies which
role, reuses that information in new contexts, and learns useful
representations. This is a source-specific architectural argument, not
a taxonomy of every computational model or a proof that one family
already explains all higher cognition
\parencite[pp.~52--64]{DoumasHummel2012HigherCognition}.

The selected presentation is the complete 2012 chapter, including its
six figures, five equations, four notes, and references. It favors
symbolic-connectionist architectures, especially dynamic binding, while
acknowledging unresolved planning, quantification, negation, and language
problems. Its model results, mathematical claims, and claims about people
have different standing. A successful simulation is not by itself evidence
that the brain implements its mechanism.

\section{What level does a model explain?}
\label{sec:p19-ccm-levels}

\begin{hbdefinition}{def:p19-ccm-levels}{Computational theory, representation and algorithm, implementation}
In the chapter's use of Marr's levels, a computational theory describes
the problem being solved and the relevant input--output relation or
normative objective. A representation-and-algorithm account specifies
what information is encoded and the procedures that operate on it.
An implementation account concerns the physical realization of those
representations and operations.

The chapter places its three process-model families primarily at the
representation-and-algorithm level. Naming nodes after neurons or
productions after cognitive operations does not establish their physical
implementation. Likewise a normative input--output standard does not
uniquely determine the process used to achieve it.
\end{hbdefinition}

The source points to normative logic/probability/choice and Bayesian
models for computational theory, and to neural models for implementation.
Those are directions within its source handbook, not replacements for
the present process-model treatment. The relevant distinction in an
evaluation is between the achieved answer, the represented information,
the operation that produced it, and the evidence for physical realization
\parencite[p.~52]{DoumasHummel2012HigherCognition}.

\section{Traditional symbolic models}
\label{sec:p19-ccm-symbolic}

\begin{hbdefinition}{def:p19-ccm-symbolic-representation}{Symbols, role bindings and recursive relational structures}
The source's symbolic system has basic representational elements and
operations for combining them into relational structures. A predicate
with ordered arguments represents a property or relation together with
its role assignments. The same predicate and object symbols can recur
in different structures, with different bindings. Recursion allows a
relation to take another relational structure as an argument.

For example, the source's gift representation assigns John the giver
role, Mary the recipient role, and a book the transferred-object role;
its heavy-book representation adds a property. Figure~5.1 instead
represents gravity causing Earth to revolve around the Sun: the
revolution proposition occupies an argument position in the causal
proposition. A labeled graph can express the same bindings if its edges
retain their roles; an unlabeled bag of nodes cannot.
\end{hbdefinition}

The chapter contrasts graph-matching approaches to analogy with ACT
production systems. Graph matching seeks corresponding substructures
between source and target; production systems apply symbolic
condition--action or condition--inference rules to a knowledge base.
They need not use the same search, control, or learning process merely
because both manipulate symbols. ACT's memory, inference, learning and
tutoring applications are surveyed claims, not newly reproduced fits
\parencite[pp.~52--53]{DoumasHummel2012HigherCognition}.

\begin{hbexample}{def:p19-ccm-transfer-rule}{The source's size, schema and fixed-object examples}
The size example supplies that a battleship is larger than a cruiser,
a cruiser than a destroyer, and a destroyer than a submarine. A rule
stating transitivity of the larger-than relation supports successive
consequences, such as battleship larger than destroyer and then larger
than submarine. The same variable-based rule can be applied to a
Rottweiler, housecat, and mouse when the corresponding premises hold.
This reuse depends on a rule applicable to those arguments; it is not
a newly learned fact about the animals' actual sizes.

Figure~5.2 also supplies a rule connecting larger-than with heavier-than.
Its presence is part of the example's model, not a universally true
physical law. Similarly, a child's modeled rule that moving things have
legs can be descriptively useful despite being false. Computational
execution and truth of the encoded rule are separate assessments.

The source's enemy-of-an-enemy stories illustrate schema induction:
two stories with corresponding relational structures can support an
abstract pattern then applied to Joe, Tim and Bill. The resulting
friendship conjecture depends on that supplied schema; it is not a
logical consequence of the meaning of enemy in every social setting.
Its passing-and-collecting examples connect Monopoly's GO square and
a border's passport stamp through an abstract pattern. The reward must
be supplied by the relevant rule or context, rather than inferred for
every possible reward from universally quantifying an unconstrained
reward variable.
\end{hbexample}

This is the symbolic strength the authors emphasize: rule variables
can take new arguments without relearning every ground instance.
Note~1 makes the qualification explicit. A rule confined to the
modeler's own neighbor's dog does not thereby generalize to all objects;
analogy can provide a different route to transfer.

The architecture has three limitations in the source. First, often
hand-coded predicates leave their acquisition unexplained. Observations
present larger-than through particular pairs; an account must explain
learning a relation that can be reused with other arguments. Second,
bare symbols do not intrinsically represent graded semantic similarity.
The token names for battleship and Rottweiler do not encode their
meaningful similarities and differences. Extra semantic machinery is
needed; the criticism concerns a bare representational scheme, not a
theorem that a symbolic system cannot store any information about
similarity. Third, brain implementation of graph matching and productions
remains a further explanatory task
\parencite[pp.~54--55 and p.~64, n.~1]{DoumasHummel2012HigherCognition}.

Semantic roles can matter separately. The source compares the agent
roles of murder and attempted murder, and the patient roles of murder
and manslaughter. Its aunt/nephew gift-giving case motivates a conjecture
about uncles and nieces through semantic similarity, not syntactic
substitution alone. Note~2 reinforces this by comparing larger with
smaller under reversed arguments: the appropriate correspondence
depends on role meaning, not merely the first printed slot. These are
theoretical examples, not a new measurement of participants' judgments.

\section{Traditional connectionist representations and processing}
\label{sec:p19-ccm-connectionist}

\begin{hbdefinition}{def:p19-ccm-network-code}{Activation, architecture, and relative localist/distributed coding}
A connectionist representation is a pattern of activation over connected
units. The source usually places activations in the interval from zero
to one. Positive connection weights are excitatory and negative weights
inhibitory. Architecture specifies which units communicate; an activation
function and update procedure specify how they do so.

Localist versus distributed coding is relative to the represented
concept. A unit interpreted as mammal may be localist for that feature
and part of a distributed dog representation. Thus the same code can
be localist for one target and distributed for another. The source
distinguishes entirely distributed vector-symbolic representations,
where a unit's interpretation depends on the whole pattern.
\end{hbdefinition}

Figure~5.3's dog, cat and shark patterns illustrate overlap among
distributed codes; overlap can make learned responses transfer between
similar inputs. This does not establish a human similarity judgment
until the feature vocabulary, coding and comparison have been justified.
Figure~5.4 compares feed-forward layers, recurrence carrying an earlier
hidden state into later processing, and auto-associative interconnection.
Their states and outputs must be defined separately. Recurrence does
not by itself guarantee convergence to a stable state or specify how
long computation takes \parencite[pp.~56--57]{DoumasHummel2012HigherCognition}.

The source credits such systems with a possible neural link, semantic
generalization, integration of learning with representation, and graded
degradation after damage. Those properties depend on the selected
architecture and learning procedure. A network is not biologically
validated simply because it has units, and a simulated lesion is not
an identified cause of a patient's impairment. The chapter surveys
learning, aging, and damage applications without supplying their whole
study records or a universal graceful-degradation theorem.

\begin{hbwarning}{def:p19-ccm-capacity-boundary}{Nonlinearity alone is not universal exact computation}
The source says a feed-forward network with three connection layers
can compute any computable mapping. As printed, the statement omits
the input domain, activation class, network size, exact-versus-approximate
criterion and learning requirements. It must not be used as a theorem
that every fixed nonlinear architecture computes every mapping.

Approximation of a class of functions, exact realization of finite
training data, and exact computation on an unbounded input domain are
different claims. Expressive capacity also does not establish that a
learning procedure finds appropriate weights or that its output captures
a person's reasoning mechanism.
\end{hbwarning}

\begin{hbtheorem}{thm:p19-ccm-continuity}{proved=annex:p19.ccm.continuity}{A precise limit on a common feed-forward interpretation}
A finite feed-forward network whose layers are affine maps followed
by coordinatewise continuous activation functions computes a continuous
real-valued map. With one rational input encoded directly as a real
number, such a network cannot output exactly zero for every input at
or below zero and exactly one for every positive rational input.
The rational sign decision is computable from a finite numerator and
positive-denominator encoding; its exact step-valued numeric map cannot
be realized under the stated continuous-network assumptions.
\end{hbtheorem}

The condition excludes a separate discontinuous threshold applied to
the network's output, discontinuous activations, and a different
discrete encoding. It is a counterexample to an unspecified blanket
guarantee, not a claim that no neural architecture can classify signs
or approximate useful functions. The source provides no complete
universal-computability theorem to reproduce at this boundary
\parencite[pp.~57--58]{DoumasHummel2012HigherCognition}.

\section{Binding and a discriminating relational task}
\label{sec:p19-ccm-binding}

The chapter identifies two informational freedoms in a symbolic
representation: which entities are represented, and how their arguments
are bound to relational roles. Merely activating loves, John and Mary
does not distinguish John loving Mary from Mary loving John. A
conjunctive code can distinguish them: John-as-lover with Mary-as-beloved
versus Mary-as-lover with John-as-beloved. The source's invented jonks/grummond/
steplock example shows that reusable symbols preserve this structural
distinction even when their meanings are unknown.

The authors argue that independent role/argument reuse is important for
analogy, schemas and rules, and criticize strictly associationist models
as capturing restricted proxies for relational reasoning. This is their
architectural position, not a mathematical result that every network
outside their favored family fails every novel-input task. In particular,
new feature combinations within a task and a new relational operation
are distinct generalization tests
\parencite[pp.~58--60]{DoumasHummel2012HigherCognition}.

\begin{hbexample}{def:p19-ccm-spatial-source}{The spatial-query example and its actual simulation record}
Figure~5.5 places a circle at row two, column two, and a square at row
three, column two in a four-by-four grid. The circle is above the square;
the relevant location is therefore row two, column two. This example
asks whether a system can use object, location and relation information
to answer a specified query.

The cited O'Reilly--Busby primary uses two objects, each encoded by two
active features among eight, arranged horizontally or vertically in
the grid. A bidirectionally connected network with 200 hidden units
has input display, object, location, relation and query representations.
Object and location layers can serve as inputs or outputs. Four query
types ask what occupies a location, where a supplied object is, or the
relation between a supplied agent object/location and the other object,
also recovering the agent location/features. The other displayed
object serves as patient. Thus the original task does contain an
agent-relative relational query; it is not merely labeling an object
above without any relational context
\parencite[pp.~4--5]{OReillyBusby2002Binding}.
\end{hbexample}

The original paper trains on 26 of the 28 possible two-feature objects,
using 200, 300 or 400 of the 416 trained-object/location combinations.
For each, it samples 4, 10, 20 or 40 patient/location combinations.
The maximum is 16,000 of 62,400 input configurations, approximately a
quarter; the minimum covers about 1.3\%. Familiar-object transfer is
assessed on a random sample of 640 untrained configurations. With
Leabra, the reported accuracy is approximately 95\% at the largest
training condition and 80\% at the 300-by-20 condition. A separate
640-item test uses the two held-out objects as agents in all locations
and 20 sampled familiar patients. Its error levels off near 15\%,
roughly 85\% accuracy, rather than establishing the same 95\% result
for genuinely held-out objects
\parencite[pp.~6--7, Figs.~4--5]{OReillyBusby2002Binding}.

This is original \emph{simulation} evidence, not a participant study.
There is no human sample or clinical population. Inputs are synthetic
and constrained to the specified feature vocabulary, spatial layouts
and query set. The text reports approximate test-error values rather
than a complete run-level uncertainty record or numerical confidence
intervals. It does not provide final weight matrices, random seeds or fitting
commands here. No simulation or parameter fit was rerun for this
chapter. The paper contrasts Leabra's learning/competition constraints
with contrastive Hebbian learning and Almeida--Pineda recurrent
backpropagation: the last failed to learn the task in the reported
attempts, and Leabra generalized better than the tested alternative.
That comparison does not identify the universally best algorithm.

The source's sharper criticism concerns a different achievement:
explicitly reusable roles and fillers supporting new inferences, not
only held-out configuration accuracy. Its falling-from-a-support example
requires a representation of which object is supported by which and
an additional rule connecting support to possible falling. The primary
simulation neither reports this rule nor tests the inference. Conversely,
its success on new objects/configurations means that the chapter's
blanket wording about only already-taught answers is too strong.
The defensible contrast is between the actual restricted generalization
shown and the further open-ended relational competence demanded
\parencite[pp.~59--60]{DoumasHummel2012HigherCognition}
\parencite[pp.~6--8]{OReillyBusby2002Binding}.

A discriminating follow-up would preserve the trained task while
separately testing unseen filler combinations, role reversals, new
relation compositions, nested propositions, and the support/falling
inference with its rule supplied. Predictions depend on the encoded
roles, vocabulary and procedure. This is a proposed test, not an
experiment already conducted; calling an item novel without specifying
which structure is novel cannot distinguish interpolation, role reuse,
and acquisition of a new rule.

\section{Symbolic-connectionist binding: products and temporal structure}
\label{sec:p19-ccm-hybrid}

\begin{hbdefinition}{def:p19-ccm-tensors}{The source's product-binding equations with explicit domains}
Let \(r\in\mathbb R^m\), \(f\in\mathbb R^n\), and
\(g\in\mathbb R^s\), with positive dimensions. Their outer-product
bindings have coordinates
\[
 (r\otimes f)_{ij}=r_i f_j,
 \qquad (r\otimes f\otimes g)_{ijk}=r_i f_j g_k.
\]
These are the chapter's equations~(1)--(2), whose displayed labels
\(rf\) and \(rfg\) denote bound tensors, not an untyped scalar
product. For a multiplace proposition, another approach sums the
separate role bindings:
\[
 T=r^{(1)}\otimes f^{(1)}+r^{(2)}\otimes f^{(2)}.
\]
This makes explicit the binding-index roles in source equation~(3).
A finite tensor can be flattened into an activation vector without
changing its listed coordinates.
\end{hbdefinition}

\begin{hbtheorem}{thm:p19-ccm-tensor-similarity}{proved=annex:p19.ccm.tensor-similarity}{Product-binding similarity and its exact boundary}
For matching finite dimensions and the ordinary Euclidean dot product,
\[
 (r^{(1)}\otimes f^{(1)})\cdot(r^{(2)}\otimes f^{(2)})
 =(r^{(1)}\cdot r^{(2)})(f^{(1)}\cdot f^{(2)}).
\]
If all four factors are nonzero, cosine similarity satisfies
\[
 \operatorname{cos}(r^{(1)}\otimes f^{(1)},r^{(2)}\otimes f^{(2)})
 =\operatorname{cos}(r^{(1)},r^{(2)})
  \operatorname{cos}(f^{(1)},f^{(2)}).
\]
The corresponding three-factor dot/norm/cosine expressions factor into
three scalar terms. A cosine involving a zero vector is undefined here.
Identical nonzero roles with orthogonal nonzero fillers give bound
cosine zero. But simultaneously reversing both factors' signs gives
the identical outer product.
\end{hbtheorem}

These are the exact identities behind source equations~(4)--(5).
They explain a whole-bound similarity problem: reusing a role need
not yield a similar tensor when its fillers differ. They do not prove
that every decoder or operation loses the role's information, that
all generalization is precluded, or that a compressed convolution
inherits every exact uncompressed identity without additional
conditions. A model's similarity procedure must be stated.

\begin{hbtheorem}{thm:p19-ccm-unbinding}{proved=annex:p19.ccm.unbinding}{A concrete counterexample to blanket tensor impossibility}
Let \(r^{(1)},\ldots,r^{(k)}\in\mathbb R^m\) be orthonormal,
with \(1\leq k\leq m\), and let fillers
\(f^{(i)}\in\mathbb R^n\) be arbitrary. For
\(T=\sum_{i=1}^k r^{(i)}\otimes f^{(i)}\), matrix decoding gives
\[
                 T^{\mathsf T}r^{(j)}=f^{(j)}\quad(1\leq j\leq k).
\]
Thus novel fillers can be recovered exactly under these stated role
and decoder assumptions. This does not explain how people learn the
roles, or establish that a particular published tensor architecture
implements this decoder.
\end{hbtheorem}

The counterexample limits the chapter's claim that product similarity
by itself rules out any generalization. It does not remove the
architectural question about semantic similarity and independently
reusable role/filler activity. Nor does it imply that arbitrary
nonorthogonal roles are perfectly decoded by the same formula
\parencite[pp.~60--61]{DoumasHummel2012HigherCognition}.

The source's favored alternative uses temporal structure as a binding
tag. Roles and their fillers fire together, with different bindings
out of phase, or a role fires immediately before its filler under
systematic asynchrony. The component representations can remain
unchanged across bindings. Describing this as additive does not mean
that an untagged sum of all active entities distinguishes every
binding: the temporal grouping is essential. Transient firing cannot
alone store a long-term relation; the chapter accordingly requires
conjunctive long-term codes and an account of their relation to dynamic
working-memory bindings.

\begin{hbexample}{def:p19-ccm-lisa}{All levels of the source's LISA example}
LISA, Learning and Inference by Schemas and Analogies, combines
distributed semantic-feature units with localist predicate/object
(PO), role-binding (RB), and proposition (P) units. In the
cup example, outside binds to cup and inside to coffee. The two
RBs are linked by a P representing cup containing coffee. Each PO
also connects to its distributed semantic features. Figure~5.6's
higher-order case adds Jim knowing that proposition: the contained
proposition occupies the known-content role rather than being replaced
by an unrelated object.

The localist levels make objects/roles, bindings and propositions
available as distinguishable entities for mapping and inference.
The different drawn shapes do not denote different neuron types;
the source explicitly says the units differ by network layer, not
by a different intrinsic update law. When a proposition is active
in working memory, PO/semantic units additionally carry the temporal
binding information. Neither the stored P nor the mere activation
of all components supplies that temporal pattern by itself.
\end{hbexample}

DORA, Discovery of Relations by Analogy, addresses acquisition.
It begins with unstructured object-feature vectors, compares objects,
learns explicit representations of shared properties, and uses time
to bind the resulting one-place predicates to new arguments. It then
combines predicates into multiplace relations. For coffee inside cup,
inside binds coffee and outside binds cup, with the container-first
convention retained in the resulting contains relation.

Note~3 says systematic asynchrony lets role and filler use the same
semantic-unit pool. Note~4 states the crucial boundary: DORA assumes
the features of relations are available in the initial representations;
it explains their extraction and explicit predication, not what those
features ultimately are. The chapter's comparison with children's
learning trajectories and its developmental/category/labeling applications
remain attributed model-and-study claims. No participant sample,
new fit or complete DORA implementation is invented from this survey
\parencite[pp.~61--64, especially nn.~3--4]{DoumasHummel2012HigherCognition}.

The authors credit symbolic-connectionist systems with structural
generalization together with distributed semantics, learning and
possible neural implementation. They nevertheless identify unresolved
executive function, goals, negation and quantification, greater
representational/processing complexity, and dependence on layer-sensitive
temporal assumptions. Their response to the nativism objection is that
this adds organizational and temporal assumptions rather than simply
hand-coding all predicates. It remains an architectural argument, not
proof that the brain uses those assumptions or that all symbolic
functions have been implemented.

\section{Evaluation, underdetermination and use}
\label{sec:p19-ccm-evaluation}

\begin{hbmethod}{def:p19-ccm-model-record}{An inspectable source-specific process-model record}
Specify the task, inputs and outputs before calling a simulation a
model of higher cognition. List the feature vocabulary, object/role
encodings, binding representation, allowable recursion, architecture,
activation/update procedure, and stopping or resource conditions.
Separate supplied predicates from learned predicates and stored bindings
from active temporal tags.

State the learning and fitting record: training examples and their
selection, held-out combinations versus held-out objects, adjusted
parameters, target/loss and learning rule, initialization, optimization
and stopping details, and the test measure. Mark unavailable details
as unavailable. A named learning algorithm is not a full reproducible
set of trained weights or hyperparameters.

Test the claimed structural achievement. Reverse roles with the same
fillers, reuse a role with unfamiliar fillers, add nested propositions,
and distinguish a new relational inference from a new spatial
configuration. Specify which rule must already be supplied. Compare
the alternatives on those discriminating cases, not only on an aggregate
accuracy that mixes them.

Report mathematical and empirical standing separately. A tensor
identity is an algebraic result; a sampled error rate is simulation
performance; human behavioral evidence needs the record in
Method~\ref{def:p19-evidence-record}; a neural mechanism needs additional
implementation evidence. Do not substitute any one for the others.
\end{hbmethod}

This is an editorial application of the chapter, not a rerun or a
universal model-selection algorithm. Its success case is a specified
spatial-model transfer result whose feature, query, training and test
boundaries are preserved. Failure cases include crediting restricted
95\% transfer as universal relational competence, claiming that a
hand-coded size predicate has been learned, and treating exact tensor
factorization or an attractive neural analogy as validation of a
cognitive mechanism.

Several models can match the same outputs with different representations
or processes. A good fit then leaves the mechanism underdetermined;
the number of observations, the parameter freedom, and whether evaluation
uses previously unseen conditions affect the inference. The fitted
response-model distinction in
Definition~\ref{def:p19-study1-response-model} illustrates why a
model-derived score is not itself an observed process.
Definition~\ref{def:p19-empirical-question} separates task, construct and measure,
and Definition~\ref{def:p19-benchmark} separates performance from its
assessment standard.

Chapter~\ref{sec:p19-analogy-and-problem-transfer} supplies a different
kind of evidence: participants' noticing, mapping and adaptation in
specified source/target problems. It cannot be replaced by executing
a graph matcher, nor does a participant's correct answer uniquely
identify the matcher or LISA as the mechanism. The present source's
Gick--Holyoak reference is the distinct 1983 schema-induction work,
not automatically the 1980 experiment series treated there.

The source concludes that a satisfactory architecture should combine
semantic content with dynamic role binding that permits reuse of
components. Symbolic-connectionist systems are offered as a promising
route, not a completed explanation of planning, language, negation,
quantification or all human reasoning. The exact mathematical results
above are fully proved in their reciprocal annex items; the remaining
architectural successes, criticisms and open questions preserve their
source-specific, dated standing
\parencite[pp.~63--66]{DoumasHummel2012HigherCognition}.

\index{computational cognitive models}
\index{symbolic representations!role binding}
\index{connectionism!relational generalization}
\index{tensor product!binding and similarity}
\index{LISA}
\index{DORA}
